\section{Case: casts in Lean 4}
\label{sec:casts}

\paragraph{Background.} A proof assistant such as Lean~4, Rocq or Isabelle/HOL checks
proofs written in a formal language. In Lean, an \emph{elaborator} turns what the author
writes into a fully explicit term, and a small \emph{kernel} checks that term against the
rules of the type theory. Every term has a type. A natural number $n : \N$ is not an
integer; where an integer is wanted, the term is ${\uparrow}n : \Z$, a different term of a
different type. Node typing (A1) holds by construction.

To stay usable, Lean gives syntactic continuity back in two ways. \emph{Typeclasses}
(\code{Add}, \code{Sub}, \code{LE}) provide one symbol and one operation per type.
\emph{Coercions} (\code{Coe} and its relatives) provide a map: where a term of type
$\alpha$ stands in a place that expects $\beta$, the elaborator inserts
${\uparrow} : \alpha \to \beta$.

\subsection{Invisible in the source, placed by the elaborator}

The inserted map is not hidden from the kernel, and Lean's goal display prints it as an
arrow. What the author does not write, and so does not see in the source, is \emph{where}
the elaborator puts it. That placement decides what the expression means:
\begin{lstlisting}
-- elaborated as ↑(2 - 3): subtract in ℕ, then cast
example : ((2 - 3 : Nat) : Int) = 0 := rfl

-- elaborated as ↑2 - ↑3: cast, then subtract in ℤ
example : ((2 : Nat) - (3 : Nat) : Int) = -1 := rfl
\end{lstlisting}
On $\N$ the symbol \code{-} is truncated subtraction, where $a - b = 0$ whenever
$a \le b$. On $\Z$ it is the group operation. The two readings agree only under a side
condition:
\[
  {\uparrow}(a - b) \;=\; {\uparrow}a - {\uparrow}b \qquad\text{provided } b \le a .
\]
Lean's library states this as \code{Int.ofNat\_sub}, with the hypothesis in its
signature. The cast from $\N$ to $\Z$ commutes with addition and multiplication
unconditionally, and with subtraction conditionally. Nothing in the symbol \code{-} or in
the arrow says so.

\subsection{Where the library typed the edge}

Mathlib \citep{mathlib2020} records the properties of the crossing in two places.

\emph{Bundled morphisms.} Beside the bare coercion, the library holds the same map as an
object whose type says what it preserves. \code{Nat.castRingHom} is the cast as a
homomorphism of semirings; an embedding $\alpha \hookrightarrow \beta$ is a map together
with a proof that it is injective; an equivalence $\alpha \simeq \beta$ carries its
inverse. \citet{baanen2022} describes the design of these bundled structures and the
classes that abstract over them.

\emph{Cast lemmas.} The tactics \code{norm\_cast} and \code{push\_cast} work from lemmas
registered for the purpose and sorted into three kinds \citep{lewis2020}. \emph{Move}
lemmas say that a cast may be moved across an operation, as in
${\uparrow}(a + b) = {\uparrow}a + {\uparrow}b$. \emph{Elim} lemmas say that a relation
between casts is a relation between the originals, as in
${\uparrow}a \le {\uparrow}b \leftrightarrow a \le b$. \emph{Squash} lemmas collapse a
chain of casts into one.

The automation reaches exactly as far as these registrations. Where a lemma is missing,
the tactic stops, and the missing lemma is the undeclared property of the edge. Practice
thus supports premise (P4) in both directions: a declared crossing can be automated, and
an undeclared one is re-established by hand.

What remains untyped is the coercion itself. An instance \code{Coe $\alpha$ $\beta$} says
that a canonical map exists. It does not say what kind of map it is or what it preserves.
That knowledge sits beside the map, lemma by lemma, in a database keyed on the
operations.

\subsection{Structure of the map and role of the crossing}
\label{sec:three-kinds}

Two things can be said about a map that crosses levels of construction in formalised
mathematics, and they are independent. One is its mathematical structure: whether it is
injective or many-to-one, and which operations it commutes with. The other is the role
in which it is used, which decides what a failure of the crossing consists in.
Table~\ref{tab:roles} gives three common roles.

\begin{table}[t]
\centering\small
\begin{tabularx}{\textwidth}{@{}P{2.3cm}P{3.2cm}P{3.0cm}L@{}}
\toprule
Role & Example & Structure of the map & What a failure consists in \\
\midrule
Extension &
  $\N \to \Z$; a subtype into its parent type &
  Injective &
  An operation has a wider meaning in the target than in the source, as subtraction
  acquires negative results. The repair is a side condition under which the two
  readings agree. \\
Identification &
  $\alpha \to \alpha/{\sim}$ &
  Many-to-one &
  A function on the source does not respect the declared relation, and so has no
  counterpart on the classes. \\
Representation &
  \code{BitVec 64} $\to \N$; a floating-point value for a real; a G\"odel number for a
  formula &
  Injective in the first example &
  Misrepresentation: the source stands for a referent in the target and does not
  correspond to it. \\
\bottomrule
\end{tabularx}
\caption{Three roles of a map between types. The first and third rows have the same
structure and different roles.}
\label{tab:roles}
\end{table}

In the first role the source is what it is and the target adds to it. In the second the
target is defined by a relation declared on the source. Here Lean does type the
crossing: \code{Quotient.lift} takes a function together with a proof that it respects
the relation, and without that proof there is no term. In the third the source is
correct or incorrect by reference to the target. For bit vectors,
\code{(x + y).toNat = (x.toNat + y.toNat) \% 2\^{}64}, so the sum of two representations
represents the sum only while nothing overflows.

The roles are not exclusive classes of maps, and a map does not determine its role. The
maps $\N \to \Z$ and \code{BitVec 64} $\to \N$ are both injective, and each commutes
with some operation only under a condition. What differs is which side sets the
standard, and that is not a property of the function. This is why the kind of a
crossing has to be declared and cannot be inferred. It is also why the kind alone is not
enough. A typed crossing needs both components of Definition~\ref{def:typed}: the kind,
which says what a failure is, and the list of what is carried, which is where the
structure of the map is recorded. Section~\ref{sec:mech-crossing} states both maps in
Lean as crossings with their side conditions.

LEM's edge taxonomy has the same shape. It has a representational type, whose failure is
misrepresentation, and a classificatory type, which holds under a declared scheme and
changes when the scheme changes, as a quotient class depends on the declared relation.

\paragraph{What the field does, and what the rule adds.} Mathlib declares what a cast
preserves, operation by operation, and its automation depends on those declarations. The
rule adds one thing: the declaration belongs to the crossing as a whole, with its kind,
and not only to pairs of a cast and an operation. Whether that would pay for itself in
Lean's design is a question for its developers, and this paper does not propose a change
to Lean.
